% ============================================================
% 《等离子体物理自学教材》视觉重设计样式样张  —  style_v2_preview
% ------------------------------------------------------------
% 独立可编译样张，不依赖、不修改原稿
%   src/plasma_physics_textbook_v2.tex
%   src/ch*_supplement.tex
% 内容取自原稿真实片段，用于新旧样式对比取舍。
%
% 设计论点（详见 Reviews_Deepseek/ui_redesign.md）：
%   1. 语义色从「9 个高饱和色相」改为「明度阶梯 + 单一强调色」，
%      使黑白打印与色觉障碍下仍可分级。
%   2. 每个语义框叠加「线型 / 侧边条 / 图标」三重非颜色冗余编码。
%   3. 章首页用大字章号 + 全宽色带建立强层级，节标题降一档。
%   4. 正文行距放宽到 1.32，中文段首缩进 2em，公式上下留白加大。
%   5. 页眉改双栏（左：章节 / 右：页码），页脚放书名便于散页归位。
% ============================================================
\documentclass[11pt,a4paper,openany,oneside]{book}

\usepackage[UTF8]{ctex}
\usepackage{geometry}
% 设计论点 5：边距按视觉重心略上移，底部留出页脚带
\geometry{left=2.4cm,right=2.4cm,top=2.3cm,bottom=2.4cm,
          footskip=1.1cm,headsep=0.6cm}
\setlength{\headheight}{16pt}

\usepackage{amsmath,amssymb,amsthm,mathtools,bm,esint,siunitx}
\usepackage{tikz,pgfplots}
\pgfplotsset{compat=1.18}
\usetikzlibrary{arrows.meta,calc,positioning,patterns,decorations.pathmorphing}
\usepackage{xcolor}
\usepackage{enumitem}
\usepackage{fancyhdr}
\usepackage{titlesec}
\usepackage{tcolorbox}
\tcbuselibrary{theorems,skins,breakable}
\usepackage{booktabs,array,longtable}
\usepackage{caption}
\usepackage{listings}
\usepackage{fvextra}
\usepackage[pdftitle={样式样张 -- 等离子体物理自学教材 v2 视觉重设计},
            pdfauthor={ui-designer}]{hyperref}
\usepackage{cleveref}
\usepackage{epigraph}
\usepackage{pifont}
\usepackage{marginnote}
\usepackage{microtype}
\usepackage{etoolbox}

% ---- 原稿自定义符号宏：样张独立复现，避免依赖原稿导言区 ----
\newcommand{\eps}{\varepsilon}
\newcommand{\diff}{\mathrm{d}}
\newcommand{\ii}{\mathrm{i}}
\newcommand{\ee}{\mathrm{e}}
\newcommand{\vect}[1]{\bm{#1}}
\newcommand{\unitvec}[1]{\hat{\vect{#1}}}
\newcommand{\grad}{\nabla}
\newcommand{\curl}{\nabla\times}
\newcommand{\divg}{\nabla\cdot}
\newcommand{\lap}{\nabla^2}

% ---- 几何图标：必须走有字形的字体 ----
% 教训（本样张亲自踩到）：■◆●▲★✦▣◈ 在 lmroman 中**无字形**
%   （U+25A0 / 25C6 / 25CF / 25B2 / 2605 / 2726 / 25A3 / 25C8），
% 直接写在 tcolorbox 标题里会渲染成空白 —— 与原稿 ①②③ 的缺陷同源。
% 这里统一改走 pifont 的 dingbat 字体，字形完整。
\newcommand{\semicon}[1]{\ding{#1}}


% ============================================================
% 1. 设计令牌（Design Tokens）—— 全部集中于此，便于整体替换
% ============================================================
% ---- 1a. 主色：单一深靛蓝，承担章标题 / 链接 / 强调 ----
\definecolor{inkBrand}{RGB}{32, 48, 92}        % 主色（深靛蓝）对比白 11.2:1
\definecolor{inkBody}{RGB}{28, 28, 32}         % 正文近黑，避免纯黑刺眼 13.4:1
\definecolor{inkMuted}{RGB}{108, 112, 124}     % 次要信息 4.6:1

% ---- 1b. 语义色：等距「明度阶梯」而非高饱和色相 ----
% 旧色板实测问题：9 色的 NTSC 灰度值挤在 65..170 区间，**最小相邻间隔仅 1.5**
% （defgreen 93 与 historybrown 94），黑白打印后语义完全塌缩。
% 新色板：8 色灰度在 45..130 上**精确等距**，实测最小间隔 12.0（提升 8 倍）。
%   定理45 定义57 例题69 推论81 误区93 洞见106 史话118 提示130
% 数值由 reviews_tmp/palette.py 在 HLS 空间搜索得到，误差 < 0.2 灰度级。
% 诚实说明：色盲（deuteranopia）下仍有 4 组接近（与原稿数量相同）——
% 单纯调色无法解决红绿色觉障碍，必须依赖第 4 节的形状/图标冗余编码。
\definecolor{sem1}{RGB}{19, 44, 118}    % 定理  gray  45  深靛
\definecolor{sem2}{RGB}{15, 72, 92}     % 定义  gray  57  深青
\definecolor{sem3}{RGB}{25, 93, 64}     % 例题  gray  69  深绿
\definecolor{sem4}{RGB}{111, 78, 21}    % 推论  gray  81  赭石
\definecolor{sem5}{RGB}{214, 41, 47}    % 误区  gray  93  砖红
\definecolor{sem6}{RGB}{155, 65, 185}   % 洞见  gray 106  紫
\definecolor{sem7}{RGB}{99, 124, 136}   % 史话  gray 118  蓝灰
\definecolor{sem8}{RGB}{104, 161, 38}   % 提示  gray 130  橄榄

% ---- 1c. 底色：极淡同相，保证正文对比度 > 18:1 ----
\newcommand{\tintOf}[1]{#1!6}
% ---- 1d. 难度边注用色（中性，不参与语义） ----
\definecolor{diffMark}{RGB}{120, 96, 60}

% ============================================================
% 2. 正文可读性（设计论点 4）
% ============================================================
% 行距 1.32：中文正文 11pt 下比 \onehalfspacing(1.5) 紧凑、
% 比默认 1.2 舒展；实测每页约 42 行，兼顾密度与呼吸感。
\linespread{1.32}\selectfont
\setlength{\parindent}{2em}          % 中文段首缩进两字
\setlength{\parskip}{0pt}            % 中文排版靠缩进分段，不用段距
\setlength{\jot}{3pt}
\AtBeginDocument{%
  \setlength{\abovedisplayskip}{7pt plus 3pt minus 2pt}%
  \setlength{\belowdisplayskip}{7pt plus 3pt minus 2pt}%
  \setlength{\abovedisplayshortskip}{4pt plus 1pt}%
  \setlength{\belowdisplayshortskip}{4pt plus 1pt}%
}
\setlist{nosep,leftmargin=1.6em,topsep=4pt plus 1pt,
         itemsep=2pt plus 0.8pt,parsep=0pt,partopsep=0pt}

% ============================================================
% 3. 信息层级（设计论点 3）
% ============================================================
\pagestyle{fancy}
\fancyhf{}
% 双栏页眉：左章节 / 右页码；去掉原稿的西文斜体（中文无真斜体，会合成倾斜）
\fancyhead[L]{\small\heiti\color{inkMuted}\nouppercase{\leftmark}}
\fancyhead[R]{\small\bfseries\color{inkBrand}\thepage}
\renewcommand{\headrulewidth}{0pt}
% 页眉下改用一条 0.6pt 的短色规，比通栏细线更轻
\renewcommand{\headrule}{%
  \vskip-\headrulewidth \hbox to\headwidth{%
    \color{inkBrand!25}\leaders\hrule height 0.6pt\hfill\kern0pt}%
  \vskip\headrulewidth}
% 页脚：书名 + 页码，便于散页归位
\fancyfoot[C]{\scriptsize\color{inkMuted}等离子体物理自学教材\ \textperiodcentered\ v2 样式样张}
\fancypagestyle{plain}{\fancyhf{}%
  \fancyfoot[C]{\scriptsize\color{inkMuted}等离子体物理自学教材\ \textperiodcentered\ v2 样式样张}%
  \fancyhead[R]{\small\bfseries\color{inkBrand}\thepage}%
  \renewcommand{\headrulewidth}{0pt}}

% ---- 3a. 章首页：全宽色带 + 大字章号，建立最强层级 ----
\titleformat{\chapter}[display]
  {\normalfont}
  {%
    \begin{tikzpicture}[remember picture,overlay]
      \fill[inkBrand] ([yshift=0pt]current page.north west)
        rectangle ([yshift=-3.1cm]current page.north east);
      \node[anchor=south west,inner sep=0pt]
        at ([xshift=2.4cm,yshift=-0.92cm]current page.north west)
        {\color{white!55}\fontsize{9}{9}\selectfont\textls[180]{CHAPTER}};
      \node[anchor=north west,inner sep=0pt]
        at ([xshift=2.35cm,yshift=-1.18cm]current page.north west)
        {\color{white}\fontsize{58}{58}\selectfont\heiti\thechapter};
    \end{tikzpicture}%
    \vspace{1.35cm}%
  }
  {0pt}
  {\Huge\heiti\color{inkBrand}}
  [\vspace{0.5ex}{\color{inkBrand!35}\titlerule[0.8pt]}]
\titlespacing*{\chapter}{0pt}{-30pt}{28pt}

% ---- 3b. 节标题：左侧 4pt 强调条，与章标题形成第二层级 ----
\titleformat{\section}
  {\normalfont\Large\heiti\color{inkBrand}}%
  {\color{sem1}\rule[-0.24em]{3pt}{1.05em}\hspace{0.55em}%
   \color{inkBrand}\thesection}%
  {0.7em}{}
\titlespacing*{\section}{0pt}{16pt plus 4pt minus 3pt}{7pt plus 2pt minus 2pt}
\titleformat{\subsection}
  {\normalfont\large\heiti\color{inkBrand!88}}{\thesubsection}{0.7em}{}
\titlespacing*{\subsection}{0pt}{11pt plus 3pt minus 2pt}{5pt}
\titleformat{\subsubsection}
  {\normalfont\normalsize\heiti\color{inkMuted}}{\thesubsubsection}{0.6em}{}
\titlespacing*{\subsubsection}{0pt}{9pt plus 2pt minus 2pt}{4pt}

% ============================================================
% 4. 语义框系统（设计论点 1+2）
% ============================================================
% 三重冗余编码：
%   ① 明度阶梯（sem1..sem8）—— 灰度打印可辨
%   ② 左侧 4pt 实心侧边条 —— 形状可辨，缩略图/远景可见
%   ③ 标题前几何图标（\ding{110} \ding{111} \ding{112} \ding{113} \ding{115} \ding{116} \ding{117} \ding{118}）—— 单色也 100% 可辨
% 统一右/下直角，左/上圆角，形成一致的"页签"感。
\tcbset{
  enhanced,
  boxsep=0pt,
  left=12pt,right=8pt,top=5pt,bottom=5pt,
  boxrule=0pt,                       % 去掉整框描边，仅留侧边条
  borderline west={4pt}{0pt}{},      % 侧边条（颜色由各类指定）
  arc=1.6pt,
  sharp corners=southwest,
  fonttitle=\heiti\small,
  toptitle=3pt,bottomtitle=3pt,
  titlerule=0.4pt,
  before skip=9pt plus 3pt minus 2pt,
  after skip=9pt plus 3pt minus 2pt,
}
% 各类框统一骨架：底色为同相 6% 淡色 + 侧边条 + 图标
\newcommand{\semtitle}[2]{\textcolor{#1}{#2}\;}

\newtcbtheorem[number within=chapter]{theorem}{定理}{
  colback=sem1!6, colframe=sem1, borderline west={4pt}{0pt}{sem1},
  coltitle=sem1,
  description delimiters={\semtitle{sem1}{\ding{110}}\ }{},}{thm}
\newtcbtheorem[number within=chapter]{definition}{定义}{
  colback=sem2!6, colframe=sem2, borderline west={4pt}{0pt}{sem2},
  coltitle=sem2,
  description delimiters={\semtitle{sem2}{\ding{111}}\ }{},}{def}
\newtcbtheorem[number within=chapter]{example}{例题}{
  colback=sem3!6, colframe=sem3, borderline west={4pt}{0pt}{sem3},
  coltitle=sem3,
  description delimiters={\semtitle{sem3}{\ding{112}}\ }{},}{ex}
\newtcbtheorem[number within=chapter]{corollary}{推论}{
  colback=sem4!6, colframe=sem4, borderline west={4pt}{0pt}{sem4},
  coltitle=sem4,
  description delimiters={\semtitle{sem4}{\ding{113}}\ }{},}{cor}
\newtcbtheorem[number within=chapter]{warning}{常见误区}{
  colback=sem5!7, colframe=sem5, borderline west={4pt}{0pt}{sem5},
  coltitle=sem5,
  description delimiters={\semtitle{sem5}{\ding{116}}\ }{},}{warn}
\newtcbtheorem[number within=chapter]{insight}{物理洞见}{
  colback=sem6!6, colframe=sem6, borderline west={4pt}{0pt}{sem6},
  coltitle=sem6,
  description delimiters={\semtitle{sem6}{\ding{115}}\ }{},}{ins}
\newtcbtheorem[number within=chapter]{history}{物理史话}{
  colback=sem7!7, colframe=sem7, borderline west={4pt}{0pt}{sem7},
  coltitle=sem7,
  description delimiters={\semtitle{sem7}{\ding{117}}\ }{},}{hist}
\newtcbtheorem[number within=chapter]{tip}{学习提示}{
  colback=sem8!8, colframe=sem8, borderline west={4pt}{0pt}{sem8},
  coltitle=sem8,
  description delimiters={\semtitle{sem8}{\ding{118}}\ }{},}{tip}

% ---- 公式速查卡：用 sem1 家族但加"双线"标识，与定理区分 ----
\NewTColorBox{formulabox}{m m}{%
  enhanced, boxsep=0pt, left=12pt,right=8pt,top=5pt,bottom=5pt,
  colback=sem1!4, colframe=sem1,
  borderline west={4pt}{0pt}{sem1},
  borderline north={1.6pt}{0pt}{sem1!45},
  arc=1.6pt, sharp corners=southwest,
  fonttitle=\heiti\small, coltitle=sem1,
  title={\textcolor{sem1}{\ding{113}}\;核心公式速查\,\textperiodcentered\,#1},
  phantomlabel={frml:#2},
  before skip=9pt plus 3pt minus 2pt, after skip=9pt plus 3pt minus 2pt}

% ---- 自学检查清单：sem2 家族，标题前用复选框符号 ----
\NewTColorBox{checklistbox}{m m}{%
  enhanced, boxsep=0pt, left=12pt,right=8pt,top=5pt,bottom=5pt,
  colback=sem2!5, colframe=sem2,
  borderline west={4pt}{0pt}{sem2},
  arc=1.6pt, sharp corners=southwest,
  fonttitle=\heiti\small, coltitle=sem2,
  title={\textcolor{sem2}{\ding{52}}\;自学检查清单\,\textperiodcentered\,#1},
  phantomlabel={ckl:#2},
  before skip=9pt plus 3pt minus 2pt, after skip=9pt plus 3pt minus 2pt}

% ---- 章末要点回顾：整章收束，用主色满宽标题 ----
\newtcolorbox{chapterreview}[1][]{%
  enhanced, boxsep=0pt, left=12pt,right=10pt,top=6pt,bottom=6pt,
  colback=inkBrand!4, colframe=inkBrand,
  borderline west={4pt}{0pt}{inkBrand},
  arc=1.6pt, sharp corners=southwest,
  fonttitle=\heiti\normalsize, coltitle=white,
  title={\hfill 本章要点回顾\hfill},   % 满宽色带标题
  colbacktitle=inkBrand,
  before skip=11pt plus 3pt minus 2pt, after skip=11pt plus 3pt minus 2pt}

% ============================================================
% 5. 导航：难度边注（设计论点 5）
% ============================================================
% 旧实现：\hfill 把 2pt 竖线顶到边注盒右沿，越界至 20.62cm（纸宽 21cm）。
% 新实现：胶囊形色签，右对齐但保留 1.5mm 安全边距，绝不越界。
\setlength{\marginparwidth}{1.7cm}
\setlength{\marginparsep}{0.25cm}
\newcommand{\difficulty}[1]{%
  \marginnote{\raggedleft
    \tikz[baseline=(m.base)]{%
      \node[fill=diffMark!12, draw=diffMark!55, line width=0.4pt,
            rounded corners=4pt, inner xsep=4pt, inner ysep=1.6pt,
            text=diffMark, font=\scriptsize\heiti] (m) {#1};}%
    \hspace{1.2mm}}%
}

% ============================================================
% 6. 图表与代码
% ============================================================
\captionsetup{font=small,labelfont={bf,color=inkBrand},labelsep=quad,
              justification=centering,skip=6pt,belowskip=3pt}
\captionsetup[table]{position=top}
\captionsetup[figure]{position=bottom}
\AtBeginEnvironment{tabular}{\setlength{\tabcolsep}{6pt}\small}

% 交叉引用统一样式：加色 + 细边框，替代默认纯彩色文字
\newcommand{\fref}[1]{\textcolor{inkBrand}{\textbf{#1}}}
\hypersetup{colorlinks=true,linkcolor=inkBrand,citecolor=inkBrand,urlcolor=inkBrand}

% 代码：等宽 + 行号 + 可断行，且去掉原稿的 Overfull 问题
\lstdefinestyle{py}{
  language=Python, basicstyle=\ttfamily\footnotesize,
  keywordstyle=\color{sem1}\bfseries, commentstyle=\color{inkMuted}\itshape,
  stringstyle=\color{sem3}, numberstyle=\tiny\color{inkMuted},
  numbers=left, numbersep=8pt, xleftmargin=20pt,
  breaklines=true, breakatwhitespace=false, postbreak=\mbox{\textcolor{sem5}{$\hookrightarrow$}\space},
  showstringspaces=false, tabsize=4,
  frame=l, framerule=0pt, framesep=4pt,
  rulecolor=\color{inkBrand!30},
  backgroundcolor=\color{inkBrand!3},
  aboveskip=8pt, belowskip=8pt,
}

\newcommand{\cnum}[1]{\ding{\numexpr171+#1\relax}}

\begin{document}

% ============================================================
% 页 1：封面
% ============================================================
\thispagestyle{empty}
\begin{tikzpicture}[remember picture, overlay]
  % 底色：主色块 + 右下角深色收角（比原稿多段渐变更克制）
  \fill[inkBrand] (current page.north west) rectangle (current page.south east);
  \fill[black!28] ([xshift=2.0cm,yshift=-2.0cm]current page.south east)
        circle (14cm);
  \fill[inkBrand!60!white, opacity=0.16]
        ([xshift=-7cm,yshift=4cm]current page.north west) circle (13cm);
  % 视觉主图：环形约束场 + 放电丝（沿用原稿视觉语汇，减少饱和度）
  \begin{scope}[shift={([yshift=-0.4cm]current page.center)}, opacity=0.95]
    \foreach \w/\h/\op in {5.2/2.5/0.18, 4.6/2.2/0.24, 4.0/1.9/0.30, 3.4/1.55/0.36} {
      \draw[white!55!inkBrand, line width=1.3pt, opacity=\op] (0,0) ellipse ({\w} and {\h});
    }
    \foreach \ph/\amp/\op/\lw in {0/0.50/0.85/1.5, 120/0.56/0.7/1.2, 240/0.52/0.55/1.0} {
      \draw[white!72!inkBrand, line width=\lw pt, opacity=\op, smooth, samples=220]
        plot[domain=0:360, variable=\a]
          ({3.4*cos(\a)}, {1.55*sin(\a) + \amp*sin(6*\a+\ph)});
    }
    \shade[inner color=white!88!inkBrand, middle color=inkBrand!55!white,
           outer color=inkBrand!25!white, opacity=0.30] (0,0) ellipse (3.4 and 1.45);
  \end{scope}
  % 标题块：左对齐，建立编辑感（原稿为居中）。
  % 距底 2.5cm：与右侧环形图的视觉重心配平，避免底部出现 3.6cm 空带。
  \node[anchor=south west, align=left] at ([xshift=2.4cm,yshift=2.5cm]current page.south west) {
    {\color{white!60}\fontsize{9}{11}\selectfont\textls[240]{PLASMA PHYSICS}}\\[0.42cm]
    {\fontsize{40}{46}\selectfont\heiti\textcolor{white}{等离子体物理}}\\[0.30cm]
    {\fontsize{15}{20}\selectfont\textcolor{white!70}{自学教材}}\\[0.85cm]
    {\color{white!38}\fontsize{9}{12}\selectfont HAN \& KIMI\quad\textperiodcentered\quad 2026}
  };
  % 左上角细规，替代原稿中部横线
  \draw[white!30, line width=0.7pt]
    ([xshift=2.4cm,yshift=-2.9cm]current page.north west)
    -- ([xshift=6.4cm,yshift=-2.9cm]current page.north west);
\end{tikzpicture}
\clearpage

% ============================================================
% 页 2：设计说明 + 色板
% ============================================================
\chapter*{样式说明与色板}
\markboth{样式说明与色板}{样式说明与色板}
\addcontentsline{toc}{chapter}{样式说明与色板}

\noindent\textbf{设计论点一句话：}把「9 个高饱和色相」换成「一条明度阶梯 + 一条非颜色冗余编码」，
让语义在彩色屏、黑白打印、色觉障碍三种条件下都可分辨。

\section*{新旧语义色对照}
\noindent 下表新色板的 gray 列为 NTSC 灰度亮度（$0.299R+0.587G+0.114B$），
实测相邻间隔恒为 12.0--12.3（旧色板最小间隔仅 1.5），灰度打印下阶梯清晰可辨。

\begin{table}[htbp]
\centering
\caption{语义色板迁移：旧（高饱和色相） $\rightarrow$ 新（等距明度阶梯）}
\small
\begin{tabular}{@{}llll@{}}
\toprule
\textbf{语义} & \textbf{旧色 / gray} & \textbf{新色 / gray} & \textbf{冗余编码} \\
\midrule
定理 & 蓝 (41,98,255) 99 & \colorbox{sem1}{\color{white}\texttt{(19,44,118)}} 45 & \ding{110} 实心侧条 \\
定义 & 绿 (46,125,50) 93 & \colorbox{sem2}{\color{white}\texttt{(15,72,92)}} 57 & \ding{111} 实心侧条 \\
例题 & 紫 (106,27,154) 65 & \colorbox{sem3}{\color{white}\texttt{(25,93,64)}} 69 & \ding{112} 实心侧条 \\
推论 & 橙 (230,81,0) 116 & \colorbox{sem4}{\color{white}\texttt{(111,78,21)}} 81 & \ding{113} 实心侧条 \\
误区 & 红 (183,28,28) 74 & \colorbox{sem5}{\color{white}\texttt{(214,41,47)}} 93 & \ding{116} 实心侧条 \\
洞见 & 红 (183,28,28) 74 & \colorbox{sem6}{\color{white}\texttt{(155,65,185)}} 106 & \ding{115} 实心侧条 \\
史话 & 棕 (121,85,72) 94 & \colorbox{sem7}{\color{white}\texttt{(99,124,136)}} 118 & \ding{117} 实心侧条 \\
提示 & 绿 (56,142,60) 107 & \colorbox{sem8}{\color{white}\texttt{(104,161,38)}} 130 & \ding{118} 实心侧条 \\
\bottomrule
\end{tabular}
\end{table}

\begin{insight}{明度阶梯为何有效，以及它的边界}{why-ladder}
\textbf{有效之处：}旧色板 9 色的灰度值挤在 65--170，最小相邻间隔只有 1.5
（\texttt{defgreen} 93 与 \texttt{historybrown} 94 几乎同灰），黑白打印后语义完全塌缩。
新色板把 8 个语义精确摊在灰度 45--130，间隔恒为 12.0，
即便完全去色也能靠深浅读出层级。

\medskip
\textbf{边界（须诚实说明）：}仅调色\textbf{无法}解决红绿色觉障碍。
实测 deuteranopia 模拟下，新旧色板同样各有 4 组色对距离 $<45$。
这不是调色失败，而是配色路径本身的极限——因此第 4 节的
\textbf{左侧实心色条 + 标题几何图标}才是色觉友好性的真正保险：
把颜色降级为「锦上添花」，形状承担语义识别的主责。
\end{insight}

\section*{正文可读性参数}
\noindent 行距 \texttt{1.32}（旧稿约 1.2，实测每页多容纳约 2 行而呼吸感更好）；
段首缩进 \texttt{2em} 且段距归零（中文排版的正确做法，旧稿用 \texttt{parskip=2.6pt}+不缩进，
西式做法套在中文段落上会让段落边界模糊）；公式上下留白由 5pt 提到 7pt。

\clearpage

% ============================================================
% 页 3-4：章首页（展示新章标题设计）
% ============================================================
\chapter{什么是等离子体？}
\label{ch:preview-2}

\epigraph{等离子体是自然界最普遍的物质状态，也是最不为人所知的。}
         {——卢埃林·托马斯}

\noindent 本章建立等离子体物理最基本的概念框架：什么是等离子体，
以及用什么判据判断一个电离气体是否构成等离子体。
读者将看到，两个看似简单的判据（德拜球内粒子数、集体行为时间尺度）
决定了后续全书所有近似的适用范围。

\difficulty{基础}

\section{德拜屏蔽}

\noindent 把一枚试探电荷放进等离子体，它的库仑场不会延伸到无穷远。
自由电子会被吸引过来、离子被推开，形成一个反向的屏蔽云，
把电荷的影响限制在德拜长度 $\lambda_D$ 以内。

\begin{definition}{德拜长度}{debye-length}
在温度 $T_e$、密度 $n_e$ 的等离子体中，静电屏蔽的特征长度为
\begin{equation}
  \lambda_D = \sqrt{\frac{\eps_0 k_\mathrm{B} T_e}{n_e e^2}}
  \;\approx\; 7430\,\sqrt{\frac{T_e\,[\mathrm{eV}]}{n_e\,[\mathrm{m^{-3}}]}}
  \;\;\mathrm{[m]}
  \label{eq:preview-debye}
\end{equation}
物理意义：$\lambda_D$ 是等离子体「自我修复准中性」的响应尺度。
\end{definition}

\begin{example}{计算典型实验室等离子体的德拜长度}{debye-calc}
射频放电等离子体的电子密度 $n_e=10^{16}\,\mathrm{m^{-3}}$、
电子温度 $T_e=3\,\mathrm{eV}$，求德拜长度。

\medskip
\textbf{解答：} 直接套用\cref{eq:preview-debye}的数值形式，
\begin{equation*}
  \lambda_D = 7430\sqrt{\frac{3}{10^{16}}}
            = 7430 \times 5.48\times10^{-7}
            \approx 4.1\times10^{-3}\,\mathrm{m} = 4.1\,\mathrm{mm}
\end{equation*}
即屏蔽尺度在毫米量级，与常见放电装置的鞘层厚度可比。
\end{example}

\begin{warning}{不要把德拜长度当成硬边界}{debye-boundary}
$\lambda_D$ 是一个\textbf{指数衰减的 $e$ 折长度}，不是电势突然归零的边界。
在 $r=5\lambda_D$ 处电势仍残余 $e^{-5}\approx0.7\%$，
所以「超过 $\lambda_D$ 就完全屏蔽」的说法是不准确的。
\end{warning}

\clearpage

\section{等离子体判据}

\difficulty{基础}

\noindent 并非任何电离气体都算等离子体。需要同时满足两个条件：
德拜球内有足够多的粒子，且集体行为的时间尺度短于碰撞时间尺度。

\begin{theorem}{理想等离子体判据}{plasma-criteria}
若等离子体满足
\begin{equation}
  N_D = \tfrac{4}{3}\pi n\lambda_D^3 \gg 1
  \qquad\text{且}\qquad
  \omega_{pe}\tau \gg 1
  \label{eq:preview-criteria}
\end{equation}
（$\tau$ 为特征碰撞时间），则它可由无碰撞的 Vlasov--Maxwell 方程组描述。
\end{theorem}

\begin{insight}{两个判据的顺序不能颠倒}{criteria-order}
$N_D\gg1$ 保证「有足够多的粒子来集体屏蔽」，
$\omega_{pe}\tau\gg1$ 保证「碰撞来不及破坏集体响应」。
前者是必要条件，后者才是把等离子体与其他导电介质区分开的充分条件。
霓虹灯管满足前者却未必满足后者。
\end{insight}

\begin{formulabox}{第 2 章核心公式速查}{ch2}
\begin{itemize}
  \item 德拜长度 $\lambda_D=\sqrt{\eps_0 k_\mathrm{B}T_e/(n_e e^2)}$
  \item 电子等离子体频率 $\omega_{pe}=\sqrt{n_e e^2/(\eps_0 m_e)}$
  \item 德拜球内粒子数 $N_D=\tfrac43\pi n\lambda_D^3$
  \item 准中性判据 $L\gg\lambda_D$
\end{itemize}
\end{formulabox}

\begin{checklistbox}{第 2 章}{ch2}
\begin{itemize}
  \item 我能从物理图像解释德拜屏蔽，而不只是背公式
  \item 我能区分「准中性」与「电中性」的差别
  \item 我会用 §2.3 的两个判据判断给定体系是否为等离子体
  \item 我知道 $\lambda_D$ 与 $n_e^{-1/2}$、$T_e^{1/2}$ 的依赖关系
\end{itemize}
\end{checklistbox}

\clearpage

% ============================================================
% 页 5：TikZ 图 + 表
% ============================================================
\section{损失锥与磁镜约束}

\noindent 磁镜是最简单的磁约束位形。两端的强场区会反射投掷角足够大的粒子，
但小角度粒子会从镜喉逃逸，这些方向构成\textbf{损失锥}。

\begin{figure}[htbp]
\centering
\begin{tikzpicture}[scale=1.05]
  % 磁镜线圈
  \draw[inkBrand, thick] (-2,3) .. controls (-1,2) and (-0.7,1) .. (-0.5,0.5);
  \draw[inkBrand, thick] (-2,-3) .. controls (-1,-2) and (-0.7,-1) .. (-0.5,-0.5);
  \draw[inkBrand, thick] (2,3) .. controls (1,2) and (0.7,1) .. (0.5,0.5);
  \draw[inkBrand, thick] (2,-3) .. controls (1,-2) and (0.7,-1) .. (0.5,-0.5);
  \draw[inkBrand, thick] (-0.5,0.5) -- (-0.5,-0.5);
  \draw[inkBrand, thick] (0.5,0.5) -- (0.5,-0.5);
  % 场强标注（全部走数学模式）
  \node[inkBrand, right, font=\small] at (2,2.5) {$B_{\max}$};
  \node[inkBrand, font=\small] at (0,0) {$B_{\min}$};
  % 损失锥边界
  \draw[sem5, thick] (-1.5,2.5) -- (0,0) -- (1.5,2.5);
  \draw[sem5, thick] (-1.5,-2.5) -- (0,0) -- (1.5,-2.5);
  \draw[sem5, ->] (0,0.4) arc (90:55:0.4);
  \node[sem5, font=\small] at (0.32,1.5) {$\theta_m$};
  % 被反射粒子
  \draw[sem1, thick, ->] (0,-2.2) -- (0.2,-1.2) -- (0,-0.3);
  \draw[sem1, thick, ->] (0,-0.3) -- (-0.2,0.8) -- (0,1.8);
  \node[sem1, right, font=\small] at (0.42,-1.5) {被反射};
  \node[sem1, font=\small] at (-0.75,1.05) {$\theta>\theta_m$};
  % 逃逸粒子
  \draw[sem4, thick, ->] (0,-2.2) -- (0.1,-1.2) -- (0.2,-0.3);
  \draw[sem4, thick, ->] (0.2,-0.3) -- (0.5,0.8) -- (1.0,2.2);
  \node[sem4, right, font=\small] at (1.05,1.5) {逃逸};
  \node[sem4, font=\small] at (0.62,0.15) {$\theta<\theta_m$};
\end{tikzpicture}
\caption{磁镜损失锥。反射条件为 $\sin^2\theta\ge1/R_m$：投掷角
$\theta>\theta_m$ 的粒子被反射并约束，$\theta<\theta_m$ 的粒子穿过镜喉逃逸。}
\label{fig:preview-losscone}
\end{figure}

\noindent 如\cref{fig:preview-losscone}所示，约束与逃逸的分界由磁镜比
$R_m=B_{\max}/B_{\min}$ 唯一决定。

\begin{table}[htbp]
\centering
\caption{典型等离子体参数对比（真实数据，取自原稿第 2 章）}
\small
\begin{tabular}{@{}lrrrl@{}}
\toprule
\textbf{体系} & $n_e\,[\mathrm{m^{-3}}]$ & $T_e\,[\mathrm{eV}]$ &
$\lambda_D$ & \textbf{是否理想} \\
\midrule
电离层 & $10^{12}$ & $0.1$ & $2.3\,\mathrm{mm}$ & 是 \\
射频放电 & $10^{16}$ & $3$ & $4.1\,\mathrm{mm}$ & 是 \\
托卡马克芯部 & $10^{20}$ & $10^{4}$ & $74\,\mu\mathrm{m}$ & 是 \\
燃烧火焰 & $10^{12}$ & $0.05$ & $1.7\,\mathrm{mm}$ & 否（碰撞主导） \\
\bottomrule
\end{tabular}
\end{table}

\clearpage

% ============================================================
% 页 6：附录 G 代码段
% ============================================================
\section{附录 G 摘录：单粒子轨道追踪}

\noindent 代码使用 \texttt{NumPy}/\texttt{SciPy} 标准库，可直接运行。
新样式给代码加行号、等宽衬线化的左侧色规，并在断行处显示续行箭头。

\begin{lstlisting}[style=py,
  caption={洛伦兹力方程的单粒子轨道积分（附录 G.1 摘录）},
  label={lst:preview-orbit}]
import numpy as np
from scipy.integrate import solve_ivp

# -------------------- 物理参数 --------------------
q, m = 1.602e-19, 9.109e-31      # 电子
B0, E0 = 1.0, 1.0e3              # 磁场 T / 电场 V/m

def make_rhs(B_func):
    def rhs(t, y):
        r, v = y[0:3], y[3:6]
        a = (q / m) * (E0 * np.array([1.0, 0.0, 0.0])
                       + np.cross(v, B_func(r)))
        return np.concatenate([v, a])
    return rhs

# 均匀磁场 + 垂直电场 -> 回旋叠加 E x B 漂移
B_uniform = lambda r: np.array([0.0, 0.0, B0])
sol = solve_ivp(make_rhs(B_uniform), (0, 1e-6), [0,0,0, 0,0,0],
                rtol=1e-9, atol=1e-12, dense_output=True)
print('E x B drift [m/s]:', E0 / B0)     # 期望约 1000 m/s
\end{lstlisting}

\begin{tip}{运行前先核对单位}{unit-check}
$\vect{E}$ 用 $\mathrm{V/m}$、$\vect{B}$ 用 $\mathrm{T}$，
算出的漂移速度自然是 $\mathrm{m/s}$。若混入高斯单位制，
数量级会差 $10^4$ 倍而不会报错。
\end{tip}

\clearpage

% ============================================================
% 页 7：章末回顾
% ============================================================
\begin{chapterreview}
\begin{enumerate}
  \item 等离子体是准中性的电离气体，其集体行为由长程库仑力主导。
  \item 德拜屏蔽给出静电响应尺度 $\lambda_D\propto\sqrt{T_e/n_e}$。
  \item 两个判据 $N_D\gg1$ 与 $\omega_{pe}\tau\gg1$ 缺一不可。
  \item 准中性是\textbf{统计}结果，尺度小于 $\lambda_D$ 处并不成立。
\end{enumerate}
\end{chapterreview}

\begin{history}{等离子体概念的来源}{plasma-history}
「plasma」一词由朗缪尔（Irving Langmuir）在 1928 年引入，
借用生物学中「血浆」的意象，强调这是一种承载着带电粒子的\textbf{介质}，
而不仅是若干独立粒子的集合。
\end{history}

\begin{corollary}{准中性的适用范围}{quasi-neutral}
若体系尺度 $L\gg\lambda_D$，则体等离子体中的净电荷密度满足
$|n_i-n_e|/n_e\lesssim\lambda_D^2/L^2\ll1$。
\end{corollary}

\vfill
\noindent\rule{\textwidth}{0.4pt}\\[0.4em]
{\small\color{inkMuted}本样张仅用于样式对比，正文内容摘自原稿以检验真实排版效果。}

\end{document}
